\documentclass[11pt,a4paper]{letter}
\usepackage[a4paper, margin=2.5cm]{geometry}
\usepackage{microtype}
\usepackage{parskip}

\signature{Marcel Bartumeu Ramentol\\
on behalf of all authors}

\address{MIT Media Lab\\
Massachusetts Institute of Technology\\
Cambridge, MA 02139, USA\\
bartumeu@media.mit.edu}

\begin{document}

\begin{letter}{Editorial Office\\
\textit{Urban Science}\\
MDPI AG, Basel, Switzerland}

\opening{Dear Editors,}

We are pleased to submit our manuscript, \textbf{``Andorra Digital Twin: A
Data-Grounded Scenario Engine for National Development Planning,''} for
consideration for publication in \textit{Urban Science}, under the Special
Issue \textit{Urban Planning for Addressing Climate Change: Technologies,
Innovations, Strategies and Modelling}.

Planning long-term national development under deep uncertainty is a critical
and underserved challenge, particularly for small countries where data are
fragmented, time series are short, and decisions span tightly coupled
socio-economic, environmental, and infrastructure systems. Existing digital
twin and scenario planning frameworks typically lack reproducible update
pipelines, cross-domain consistency, and transparent causal assumptions---gaps
that are especially consequential in data-scarce contexts.

This paper addresses these gaps by introducing the Andorra Clone Scenario
Engine, a harmonized, data-grounded framework that produces structured
long-horizon trajectories (to 2049) under four contrasting development
pathways. A central methodological contribution is the inversion of the
conventional population--economy causal chain: rather than treating population
as an exogenous driver, our model derives it endogenously from GDP growth via a
calibrated labour-immigration elasticity grounded in IMF and World Bank data
for Andorra---a micro-state where 55\% of residents are foreign-born labour
migrants. Three core parameters are calibrated against the 2010--2024
historical record using Nelder--Mead optimisation (RMSE\,=\,1.77\%), and two
structural nonlinearities---a housing-affordability migration brake and a
logistic tourism carrying-capacity ceiling---are explicitly modelled. The
framework also integrates a tangible, human-in-the-loop physical interface,
consistent with the MIT Media Lab City Science vision of transparent,
community-facing decision support.

We believe this manuscript is well suited to \textit{Urban Science} and to the
Special Issue for three reasons. First, it demonstrates a replicable
methodology for scenario-based territorial planning in data-scarce, climate-
and resource-constrained environments. Second, it makes climate and
sustainability pressures central to the scenario design---directly engaging
carrying-capacity limits, natural coverage, CO$_2$ trajectories, and water
security. Third, it bridges digital twin modelling, human-in-the-loop
interaction, and national-scale governance, spanning the technological,
strategic, and modelling dimensions called for by the Special Issue.

This manuscript has not been submitted previously to any MDPI journal or any
other publication.

We confirm that neither the manuscript nor any parts of its content are
currently under consideration for publication with or published in another
journal. All authors have approved the manuscript and agree with its submission
to \textit{Urban Science}.

We look forward to the editors' consideration.

\closing{Sincerely,}

\end{letter}
\end{document}
